% =====================================================================
% Clase -- Física 10° -- Trimestre I -- Semana 04
% Sesión 007 (jue 24 sep., 7:50, 100 min) y sesión 008 (vie 25 sep., 7:50, 100 min)
% Tema: Vectores
% Subtemas: Representación, componentes y suma gráfica / Suma por componentes y fuerza neta
% Hilo conductor del trimestre I: «De la inercia a los Principia»
% Tema 2 de la guía: Vectores (tema:vectores)
% Material del docente (no se entrega a los estudiantes).
% =====================================================================
\documentclass[11pt]{article}
\makeatletter\def\input@path{{../../../../../plantillas/estilo/}}\makeatother
\usepackage{mmcantillo}
\guiadelestudiante{../../guia-didactica/guia-periodo-I-leyes-newton}

\tipodocumento{Clase}
\asignatura{Física}
\titulo{Semana 4: cuando la dirección importa}
\grado{10°}
\periodo{I}

% DBA: naturales grado 10 · DBA 1 — Comprende que el reposo o el movimiento rectilíneo
% uniforme se presentan cuando las fuerzas aplicadas sobre el sistema se anulan entre ellas,
% y que en presencia de fuerzas resultantes no nulas se producen cambios de velocidad.
% Evidencia de esta semana: la herramienta para sumar fuerzas. Sin vectores no se puede
% decidir si «se anulan entre ellas», que es literalmente lo que dice el DBA.

\begin{document}

\encabezado*

\begin{center}
{\renewcommand{\arraystretch}{1.3}
\begin{tabularx}{\linewidth}{@{}l l l X l@{}}
  \toprule
  \textbf{Sesión} & \textbf{Fecha} & \textbf{Min} & \textbf{Subtema} & \textbf{Entregables}\\
  \midrule
  007 & jue 24 sep., 7:50 & 100 & Representación, componentes y suma gráfica & Taller\\
  008 & vie 25 sep., 7:50 & 100 & Suma por componentes y fuerza neta & Se deja tarea\\
  \bottomrule
\end{tabularx}}
\end{center}

Segunda semana del trimestre. La medición nos dio números con unidades; esta semana esos
números aprenden a apuntar. Es la herramienta que le faltaba a Aristóteles y que Newton da por
supuesta en los \emph{Principia}: para saber si un cuerpo está en equilibrio no basta con
cuánto empuja cada fuerza, hay que saber \emph{hacia dónde}. Referencia: \refguia{tema:vectores}.

% ---------------------------------------------------------------------
\seccion{Sesión 007 — jueves 24 de septiembre · 7:50 · 100 min}

\textbf{Objetivo.} Representar un vector, descomponerlo en sus componentes y sumar dos
vectores gráficamente.

\textbf{Materiales.} Tablero; regla, transportador y papel cuadriculado; calculadora
científica (en modo grados); una cuerda o un lazo para la demostración del arranque.

\begin{center}
{\renewcommand{\arraystretch}{1.3}
\begin{tabularx}{\linewidth}{@{}l l X@{}}
  \toprule
  \textbf{Momento} & \textbf{Min} & \textbf{Actividad}\\
  \midrule
  Demostración & 0--12 & Dos estudiantes halan una cuerda atada a una silla: primero en el
    mismo sentido, después en sentidos opuestos, después en ángulo. Misma fuerza de cada
    uno, tres resultados distintos. La pregunta del día: \emph{¿cuánto suman \qty{30}{N} y
    \qty{40}{N}?} Respuesta honesta: depende.\\
  Escalar y vectorial & 12--25 & Qué magnitudes necesitan dirección (fuerza, velocidad,
    desplazamiento) y cuáles no (masa, tiempo, temperatura). Notación: $\vec{F}$, flecha,
    módulo, dirección y sentido. Escala en el dibujo: \qty{1}{cm} por cada \qty{10}{N}.\\
  Componentes & 25--50 & Proyectar sobre los ejes: $F_x = F\cos\theta$, $F_y = F\sen\theta$.
    Ejemplo guiado: \qty{50}{N} a $37^\circ$ da $\qty{39,9}{N}$ y $\qty{30,1}{N}$.
    Y al revés, de componentes a magnitud y ángulo, con Pitágoras y la tangente inversa.
    \textbf{Ojo al cuadrante} (ver la nota).\\
  Suma gráfica & 50--70 & Método del paralelogramo y método del polígono (cabeza con cola),
    con regla y transportador. Comprobar el caso \qty{30}{N} este $+$ \qty{40}{N} norte:
    el dibujo da \qty{50}{N} y la cuenta también.\\
  Taller & 70--95 & Taller individual (5 puntos).\\
  Cierre & 95--100 & Recoger. Pregunta abierta: «¿y si son tres fuerzas, o cinco?
    ¿Dibujamos todo?» Mañana, el método que no necesita dibujo.\\
  \bottomrule
\end{tabularx}}
\end{center}

% ---------------------------------------------------------------------
\seccion{Sesión 008 — viernes 25 de septiembre · 7:50 · 100 min}

\textbf{Objetivo.} Sumar cualquier número de vectores por componentes y calcular la fuerza
neta sobre un cuerpo.

\begin{center}
{\renewcommand{\arraystretch}{1.3}
\begin{tabularx}{\linewidth}{@{}l l X@{}}
  \toprule
  \textbf{Momento} & \textbf{Min} & \textbf{Actividad}\\
  \midrule
  Revisión & 0--12 & Revisar el taller, sobre todo el último punto (entre qué valores puede
    estar la suma de \qty{30}{N} y \qty{40}{N}). Dejar claro que el rango va de \qty{10}{N}
    a \qty{70}{N}.\\
  El método & 12--40 & Sumar por componentes: se descompone cada vector, se suman todas
    las $x$ por un lado y todas las $y$ por otro, y al final se recompone. Tabla en el
    tablero con una columna por vector. Ejemplo con tres fuerzas
    (\qty{20}{N} a $0^\circ$, \qty{30}{N} a $90^\circ$, \qty{10}{N} a $180^\circ$):
    neta $\approx \qty{31,6}{N}$ a $71{,}6^\circ$.\\
  Fuerza neta & 40--60 & Nombrar $\vec{F}_{\text{neta}} = \sum \vec{F}$ y conectarlo con el
    DBA: si la neta es cero, el cuerpo sigue como estaba; si no, cambia su velocidad. Es el
    anuncio literal de la primera ley, que empieza en la sesión 012.\\
  Equilibrio & 60--80 & Un caso con neta cero: dos fuerzas iguales y opuestas, o tres que
    se cierran formando un triángulo. Que lo comprueben por componentes y vean los ceros.\\
  Cierre & 80--100 & Dejar la tarea. Anuncio de la semana 05: repaso de cinemática — y
    aviso de que MUA y caída libre son tema nuevo para este grupo, no repaso.\\
  \bottomrule
\end{tabularx}}
\end{center}

\begin{nota}
\textbf{El error del cuadrante.} Con componentes $(-6, 8)$ la calculadora, si se teclea
$\tg^{-1}(8 \div -6)$, devuelve $-53{,}1^\circ$ — que apunta justo al lado contrario del
vector real, que está a $126{,}9^\circ$. La calculadora no sabe en qué cuadrante estamos.
Regla de la clase: \textbf{primero se dibuja el vector a mano alzada, después se calcula}; si
el número no concuerda con el dibujo, el número está mal.
\end{nota}

\begin{nota}
\textbf{Si el grupo va lento} (son 100 minutos, pero el tema es denso): la suma gráfica del
jueves se puede recortar al método del paralelogramo, dejando el del polígono para cuando
aparezcan tres o más fuerzas el viernes. Lo que no conviene recortar es la demostración de la
cuerda: es lo que hace obvio por qué \qty{30}{N} y \qty{40}{N} no siempre son \qty{70}{N}.
\end{nota}

\begin{nota}
Los seis ejercicios de esta semana son nuevos en el banco
(\texttt{recursos/banco/fisica/vectores-10.py}): no había ninguno de vectores, aunque
\texttt{cinematica-10} y \texttt{leyes-newton-10} los usan todo el tiempo.
\end{nota}

\end{document}
